\section{A short-arc range for mixed inert C and D words}

% BEGIN M82R1
\paragraph{M82R1. Cuts with one prescribed thin letter on every mixed orbit.}
Fix a guarded inert C or D place and a good coefficient test in
I82W1--I82W5. Suppose its cyclic word has both a thick and a
nonthick letter. Set $q=p-1$, $h=B+H+S_0$ and take
\[
    j=1,\qquad G'=g=\max\{1,H,S_0,
                         \max_{a\ {\rm nonthick}}(s_a+S_0)\}.
\]
Keep the up or down sign and size of each thick letter and mark
all other letters thin. Their jumps are in $[-g,g]$ by I82W2.
On thick letters the translated cloud has width $S_0$, the
central width is $H$, and $s_a+S_0\le h$. In particular
$g\le h$ and $s_a>g$ at each thick position for a sufficiently
late member, by I82W3. These statements hold for both tested
lattices, the kernel by strictness, and the crystalline extension
of the target by the trivial representation used in I82B5.

Cut at every position immediately preceding a thin letter. The
set is nonempty. In the backward arc ending at $c$ its length
is the cyclic distance to the previous cut, interpreted as a full
period for a single cut. I82W5 excludes three consecutive thick
letters on a mixed orbit. Thus $1\le\ell_c\le\min(3,n_w)$,
the earliest letter is thin, and its at most two thick letters
have backward offsets zero or one. This verifies consecutive
cut, no-interior-cut, thin-after-cut and two-mark requirements,
including a cut followed by another cut. Let $\ell$ be the
greatest of these lengths.

Use the two original upper and lower height sequences of I82W2
in $rr_{\rm par}$; their sums are at most $2h$, and at the
letter immediately after each cut at most $2g$.
Take $M,\sigma,R=rr_{\rm par},\kappa$ and $a_0$ exactly as in
I82B2 for the above $j,g,h,\ell$. Its estimates give
\[
 a_0=2\sigma+R+2g\ell
     \le 8h/q^2+11g+2\log_p h+12 .
\]
Indeed $\sigma\le2h/q^2+2g+3$ and
$R\le g+4h/q^2+2\log_p h+6$. Uniformly on the late family,
$g/B\to0$ since $g\le2b'$ and $b'/B\to0$, while $h/B\to1$
and $\log_p h/B\to0$. As $q\ge10$, the limsup of $2a_0/B$
is at most $16/100$. Thus $2a_0<B$ eventually on every such
word. For two thick sizes $s_0,s_1\in[B,B+H]$ one has
$s_0<p s_1$ and $2s_0\le p s_1$ eventually. This proves
the remaining order-sensitive guards.

The local construction of B82R1--B82R3 through its first
conclusion now applies with these parameters. In particular the
powers at the first thin letter use $j=1$, all other marked
positions have $s<\ell_c-j$, and the integral truncation of
I82B4 allows $j=1$. The two chosen upper heights at the cut
give R82W4, I82B5 decomposes the target, and the integral
constant projection and both polynomial bounds use only the
verified size inequalities and the two-mark condition.
Consequently every finite target class lifts at exponent
\[
          E\le32h/q^2+122g\ell+14\log_p h
             \le32h/q^2+366g+14\log_p h .
\]
This comes from the local cut construction before comparison
with any family row allowance.
% END M82R1

% BEGIN M82R2
\paragraph{M82R2. Absorbing the mixed-orbit exponent in its frozen row.}
Use $F_w$ and $d=n_w$ as in N82I2. For any tested mixed
reading the local row colength is at most
$(d+1)r\,(32h/q^2+366g+14\log_p h)$ by M82R1 and I82B1.
It is therefore bounded by the stored $\mathrm{lenRowCD}(w)$ if
\[
 32h/q^2+366g+14\log_p h
              \le368h_B/p^2+2082 F_w .
\]
This uses exactly the high and flat terms of the lower endpoint
of the native interval; the multiplicity term is nonnegative.
In particular a sufficient condition on this orbit is
\[
                         p^2g\le \tfrac89 h_B .
\]
Indeed $h\le h_B$, $p/q\le11/10$ and
$32(11/10)^2+366(8/9)<368$. On the given family
$b=\exp(\mathrm{large}_2)$ and $h_B\le2b$ eventually by
the parameter bounds. Thus $\log_p h\le3+\mathrm{large}_2$
using $p\ge11$. As $\mathrm{large}_2\ge1$, $14\log_p h$
costs at most $56\mathrm{large}_2\le2082F_w$,
independently of band membership. The $g$ in these tests uses
the nonthick sizes on the actual orbit. To bound a place
uniformly one can impose the condition on every good reading
there, or use the single upper bound
$\max\{1,H,W_0+S_0\}$ for all of them. The tests also allow
different places of one index to satisfy different local ranges,
as in G82I2.

Fix any real $C\ge0$ before taking the cutoff in $x$. Then the
sufficient condition above holds at every mixed reading with
$p\le x^C$, for all members after one cutoff depending on $C$
and the family. Indeed $g\le2b'$, $h_B\ge b$, and I82W1 gives
\[
     \frac{p^2g}{h_B}\le2x^{2C+1}\exp(-x^{\epsilon/3})
                      \longrightarrow0 .
\]
This covers every mixed orbit of the small-order inert row D
after a single family cutoff: that row has $p<\ell_A$, where
$\ell_A=k x^\epsilon/(m\,\mathrm{de}_2)$. Since $k\le x^5$
by A78P2, $m\ge2$ and $\mathrm{de}_2\ge x^{-\epsilon}$
by A78P3, one has $\ell_A\le x^{5+2\epsilon}<x^6$.
Use $C=6$. The same argument covers every mixed medium inert
reading through any fixed power of $x$. At an index whose
Frobenius is different from complex conjugation, I82W5 rules
out all-thick words, and N82I1 covers its entirely nonthick
places. Therefore the global native inequality holds throughout
row D at each such index, by G82I2.
% END M82R2

% BEGIN M82R3
\paragraph{M82R3. The three-term allowance on words of every shape.}
At any guarded inert C or D place put
$E_p=1+h_B/p+\log_p h_B$. The original $\mathrm{coreM}$
is $dm E_p$. By N82I2's coarse Wach lift and the generator
bound of I82B1, the normalized row colength is at most
$(d+1)r(25h_B/p+3)$ with no condition on the word. It therefore
satisfies the stored row bound whenever
\[
  (d+1)r(25h_B/p+3)
  \le (d+1)r\bigl(368h_B/p^2+2082F_w\bigr)
                            +22365\,dm E_p.
\]
All three summands on the right are the prescribed lower
endpoint of the rational interval. In particular this covers
every prime and every coefficient test at the place under
$25(d+1)r\le22365\,dm$, since $h_B\ge1$ and
$25h_B/p+3\le25E_p$. The condition $r\le596m$ suffices
simultaneously at every place of an inert index, since the
letter orbit has even positive length by I82W5, hence $d\ge2$.
These comparisons include all-thick words. They do not impose
this optional rank inequality on the family.
% END M82R3
